\section{The rule}

A buyer agent at one company has read a refund case and confirmed the customer,
and it asks a vendor agent at a second company to stage the refund against the
vendor's payment ledger. The vendor's software must decide, before it touches
that ledger, whether to obey. Today that decision rests on an API key and a
transport signature, which establish who is calling and nothing about whether
this call, with these arguments, under the agreement these
two companies signed, was already authorized on the vendor's side, already
performed, or rests on a delegation revoked four minutes ago. Afterward each
company holds a log it wrote about itself, and no third party can tell whether
the two describe the same event. Having the sender prove more makes the
receiver's security a function of the sender's correctness, and a shared ledger
or broker makes that substrate a participant in every call.

The rule here avoids both. \dfn{Authority is minted and verified only inside the
kernel that dispatches, against roots that kernel activated itself, and nothing
carried in a request can add a root. Only signed, content-addressed evidence
crosses an organization boundary.} A signature proves who signed some bytes. It
cannot decide whether the receiver should act on them, because that decision
depends on facts that exist only at the receiver: which keys it pinned, which
agreement it activated and when, whether this continuation has been spent.

That is a placement argument; it composes with the rest of the stack rather
than replacing it \cite{saltzer1984}. Alignment training, transcript monitors,
and the sender's own policy engine change how often the kernel says no; none can
know that the callee's budget is exhausted or that this argument hash was
already dispatched. The receiver decides.

The receipt is the one object this paper asks a reader to hold. Locally a
\dfn{receipt} is a kernel-signed record of one admission decision whose identity
is the hash of its own content; across a boundary the same object is
\dfn{co-signed}, a statement whose single subject is the digest of a receipt the
receiving kernel already holds, signed by exactly two keys the receiver pinned
before the request arrived. A cross-organization call is that admission
operation with the peer's receipt and co-signature supplied as evidence, never
as authority, which makes the statement a pointer rather than a bearer
credential: holding it grants nothing, because everything it names must already
be in the receiver's stores, and no deviation is one the sender can repair by
being more persuasive. Agent protocols define how a tool is called and how
agents talk \cite{mcp,a2a}; neither defines the record of what was admitted. We
call the mediating component the kernel; the implementation measured here, Chio,
is userspace and not an operating-system kernel.
